\documentclass[10pt,a4paper]{article}
\usepackage{fontspec,amsmath,amssymb,graphicx}
\usepackage[margin=18mm]{geometry}
\setmainfont{Times New Roman}
\setlength{\parindent}{0pt}
\setlength{\parskip}{5pt}
\allowdisplaybreaks
\sloppy
\emergencystretch=2em
\begin{document}
\section*{BT\_5\_5\_3\_09 — Lời giải bổ sung}
Bochner: $\Delta|\nabla u|^2=2\sum_{i,j}u_{ij}^2+2\nabla u\cdot\nabla\Delta u=2\sum u_{ij}^2\ge0$. Nguyên lý cực đại cho hàm dưới điều hòa suy supremum ở biên. Nếu giá trị cực đại đạt trong miền, nguyên lý mạnh buộc $|\nabla u|^2$ hằng và $D^2u=0$, tức $u$ affine; khi ấy cực đại cũng đạt trong miền. Vì thế chữ ``chỉ'' trong đề cần ngoại lệ affine/hằng.\end{document}